\chapter{Flashing the Pi, headless}
\label{ch:flashing}

\chapterbrief
  {Write an operating system to the SD card, pre-configured so the Pi joins your Wi-Fi and accepts SSH on first boot.}
  {Headless setup --- no monitor, no keyboard, ever.}
  {A Pi that boots and appears on your network.}

\section{Headless means no screen}

You are never going to plug a monitor into this Pi. Instead you configure it
\emph{before} first boot, and from then on you reach it over the network.

\begin{analogy}{Setting up a phone before you unbox it}
Imagine you could type your Wi-Fi password and set your account name into a
phone \emph{before} it first powered on --- so that the moment it boots it is
already on your network and already yours.

That is what the Raspberry Pi Imager's settings do. The alternative is finding a
spare HDMI cable, a keyboard, and somewhere to put them.
\end{analogy}

\section{Raspberry Pi Imager}

Download it from \code{raspberrypi.com/software} and install it. It runs on
macOS, Windows and Linux.

Insert your SD card, open Imager, and make three choices:

\begin{description}[leftmargin=0pt, style=nextline]
  \item[Device] Your Pi model --- e.g.\ \emph{Raspberry Pi Zero 2 W}.
  \item[Operating System] \textbf{Raspberry Pi OS Lite (64-bit)}, under
        \emph{Raspberry Pi OS (other)}.
  \item[Storage] Your SD card. \textbf{Read this twice.}
\end{description}

\begin{trap}{Imager will happily erase the wrong disk}
It writes to whatever you select and does not undo. Before clicking, confirm the
\emph{size} matches your SD card --- a 32\,GB entry is your card, a 1\,TB entry
is your laptop.

Unplug other external drives while you do this. People have lost backups here.
\end{trap}

\begin{concept}{Why ``Lite'', and why 64-bit}
\textbf{Lite} means no desktop environment. You do not need one --- you will use
SSH --- and on a 512\,MB Zero 2 W a desktop consumes memory your builds need.

\textbf{64-bit} matters because the Java 21 packages are better supported there,
and the JVM performs measurably better.

(Volume One's kiosk chapter does need a graphical desktop. You will install just
the pieces required at that point, rather than carrying a full desktop through
every chapter before it.)
\end{concept}

\section{The settings screen --- the important part}

Before writing, click the \textbf{gear icon} (or answer \emph{``Would you like to
apply OS customisation settings?''} with \textbf{Edit Settings}).

\subsection*{General tab}

\begin{center}\small
\begin{tabular}{@{}L{0.30\linewidth}L{0.62\linewidth}@{}}
\toprule
\thead{Field} & \thead{Set it to} \\
\midrule
Set hostname & \code{brewer} \newline
  {\small This is how you will address the Pi: \code{brewer.local}. Pick
   something short you will not mistype.} \\
Set username and password & Username \code{martin} (or yours), and a real
  password. \newline
  {\small You will only type this password a handful of times before
   Chapter~\ref{ch:sshkeys} replaces it with a key.} \\
Configure wireless LAN & Your Wi-Fi name (SSID) and password, and your
  \textbf{Wireless LAN country} --- required, or Wi-Fi stays off. \\
Set locale settings & Your timezone and keyboard layout. \\
\bottomrule
\end{tabular}
\end{center}

\begin{trap}{The Zero 2 W is 2.4\,GHz only}
It cannot see a 5\,GHz network. If your router presents both bands under
different names, enter the \textbf{2.4\,GHz} one.

If both bands share one name, it usually still works --- but ``the Pi never
appears on the network'' with correct-looking settings is very often this.
\end{trap}

\subsection*{Services tab}

\textbf{Enable SSH}, and choose \textbf{Use password authentication} for now.
You will switch to keys in Chapter~\ref{ch:sshkeys}.

\begin{concept}{What SSH is}
\emph{Secure Shell}: a program that gives you a terminal on another machine over
the network, with everything encrypted in transit.

When you type \code{ssh martin@brewer.local}, you get a shell prompt on the Pi.
Commands you type run \emph{there}. It is the single most important tool in
embedded Linux work, and Chapter~\ref{ch:sshkeys} is entirely about making it
pleasant.
\end{concept}

\section{Write it}

Click \textbf{Save}, then \textbf{Write}, and confirm. It takes a few minutes ---
it writes, then verifies.

When it finishes, macOS may complain that ``the disk you inserted was not
readable''. \textbf{That is normal}: the card now has a Linux filesystem your
laptop cannot read. Eject it and ignore the warning.

\section{First boot}

\begin{enumerate}
  \item Put the SD card in the Pi.
  \item On a Zero 2 W, plug power into the port marked \textbf{PWR IN} --- the
        one nearer the corner. The other is USB data and will not power it.
  \item Wait. \textbf{First boot takes two to three minutes} --- it resizes the
        filesystem and reboots itself once.
\end{enumerate}

The green LED tells you what is happening: irregular flickering is disk
activity, which is good. A steady green light with no flicker for several
minutes usually means it has not booted.

\begin{checkpoint}{The Pi is on the network}
From your laptop:

\begin{shell}
ping brewer.local
\end{shell}

\begin{trace}
PING brewer.local (192.168.1.42): 56 data bytes
64 bytes from 192.168.1.42: icmp_seq=0 ttl=64 time=4.331 ms
\end{trace}

Ctrl-C to stop. Note the IP address --- you may need it in the next chapter.

\textbf{If it does not answer}, wait a full three minutes and try again before
troubleshooting. If it still fails, go to Chapter~\ref{ch:contact}, which is
entirely about finding a Pi that will not answer to its name.
\end{checkpoint}
